\SecondArticle{6}
Wir betrachten das Produkt aus $m(m-1)$ Faktoren
\[
\begin{aligned}
&(a-b)(a-c)(a-d)\, .\, .\, .\, .\, .\\
\times{}&(b-a)(b-c)(b-d)\, .\, .\, .\, .\, .\\
\times{}&(c-a)(c-b)(c-d)\, .\, .\, .\, .\, .\\
\times{}&(d-a)(d-b)(d-c)\, .\, .\, .\, .\, .\\
&\text{usw.}
\end{aligned}
\]
das wir mit $\pi$ bezeichnen. Da es die Unbestimmten $a,b,c$ usw. symmetrisch enthält, nehmen wir an, es sei auf die Form eines Polynoms in $\lambda',\lambda'',\lambda'''$ usw. gebracht. Dieses Polynom gehe in $p$ über, wenn anstelle von $\lambda',\lambda'',\lambda'''$ usw. jeweils $l',l'',l'''$ usw. eingesetzt werden. Danach nennen wir $p$ die \textit{Determinante} des Polynoms
\[
y=x^m-l'x^{m-1}+l''x^{m-2}-l'''x^{m-3}+\text{usw.}
\]
\SourceLeaf{046}
So haben wir zum Beispiel für $m=2$
\[
p=-{l'}^2+4l''
\]
Ebenso findet man für $m=3$
\[
p=-{l'}^2{l''}^2+4{l'}^3l'''+4{l''}^3-18l'l''l'''+27{l'''}^2
\]
Die Determinante des Polynoms $y$ ist also ein Polynom in den Koeffizienten $l',l'',l'''$ usw., das durch die Substitutionen $l'=\lambda',l''=\lambda'',l'''=\lambda'''$ usw. in das Produkt aus allen Differenzen zwischen je zweien der Größen $a,b,c$ usw. übergeht. Im Fall $m=1$, wenn also nur eine einzige Unbestimmte $a$ und daher überhaupt keine Differenz vorhanden ist, empfiehlt es sich, die Zahl 1 selbst als Determinante des Polynoms $y$ anzunehmen.

\Needspace{8\baselineskip}Bei der Festlegung des Begriffs der Determinante mussten die Koeffizienten des Polynoms $y$ als unbestimmte Größen betrachtet werden. Die Determinante eines Polynoms mit bestimmten Koeffizienten,
\[
Y=x^m-L'x^{m-1}+L''x^{m-2}-L'''x^{m-3}+\text{usw.}
\]
ist eine bestimmte Zahl $P$, nämlich der Wert des Polynoms $p$ für $l'=L',l''=L''$, $l'''=L'''$ usw. Nehmen wir also an, $Y$ lasse sich in lineare Faktoren zerlegen,
\[
Y=(x-A)(x-B)(x-C)\, .\, .\, .\, .\, .
\]
oder $Y$ entstehe aus
\[
\upsilon=(x-a)(x-b)(x-c)\, .\, .\, .\, .\, .
\]
indem man $a=A,b=B,c=C$ usw. setzt, sodass dieselben Substitutionen $\lambda',\lambda'',\lambda'''$ usw. jeweils in $L',L'',L'''$ usw. überführen, so ist $P$ offenbar gleich dem Produkt aus den Faktoren
\[
\begin{aligned}
&(A-B)(A-C)(A-D)\, .\, .\, .\, .\, .\\
\times{}&(B-A)(B-C)(B-D)\, .\, .\, .\, .\, .\\
\times{}&(C-A)(C-B)(C-D)\, .\, .\, .\, .\, .\\
\times{}&(D-A)(D-B)(D-C)\, .\, .\, .\, .\, .\\
&\text{usw.}
\end{aligned}
\]
Ist also $P=0$, so müssen sich unter den Größen $A,B,C$ usw. wenigstens zwei gleiche befinden; ist dagegen nicht $P=0$, so sind sämtliche Größen $A,B,C$ usw. notwendig ungleich. Nun bemerken wir: Setzt man $\frac{\mathrm dY}{\mathrm dx}=Y'$, oder

\SourceLeaf{047}
\[
\begin{gathered}Y'=mx^{m-1}-(m-1)L'x^{m-2}+(m-2)L''x^{m-3}\\-(m-3)L'''x^{m-4}+\text{usw.}\end{gathered}
\]
so gilt
\[
\begin{aligned}
Y'={}&(x-B)(x-C)(x-D)\, .\, .\, .\, .\, .\\
&+(x-A)(x-C)(x-D)\, .\, .\, .\, .\, .\\
&+(x-A)(x-B)(x-D)\, .\, .\, .\, .\, .\\
&+(x-A)(x-B)(x-C)\, .\, .\, .\, .\, .\\
&+\text{usw.}
\end{aligned}
\]
\Needspace{16\baselineskip}Sind also zwei der Größen $A,B,C$ usw. gleich, etwa $A=B$, so ist $Y'$ durch $x-A$ teilbar; mit anderen Worten, $Y$ und $Y'$ enthalten den gemeinsamen Teiler $x-A$. Nimmt man umgekehrt an, $Y'$ habe mit $Y$ irgendeinen gemeinsamen Teiler, so muss $Y'$ einen der linearen Faktoren $x-A,x-B,x-C$ usw. enthalten, etwa den ersten, $x-A$. Dies kann offenbar nur geschehen, wenn $A$ einer der übrigen Größen $B,C,D$ usw. gleich ist. Aus alledem erhalten wir also zwei \textsc{Sätze}:

I. \textit{Ist die Determinante des Polynoms $Y$ $=0$, so hat $Y$ gewiss mit $Y'$ einen gemeinsamen Teiler; haben also $Y$ und $Y'$ keinen gemeinsamen Teiler, so kann die Determinante des Polynoms $Y$ nicht $=0$ sein.}

II. \textit{Ist die Determinante des Polynoms $Y$ nicht $=0$, so können $Y$ und $Y'$ gewiss keinen gemeinsamen Teiler haben; oder, wenn $Y$ und $Y'$ einen gemeinsamen Teiler haben, muss die Determinante des Polynoms $Y$ notwendig $=0$ sein.}

\SecondArticle{7}
Doch ist sorgfältig zu beachten, dass die ganze Kraft dieses sehr einfachen Beweises auf der Annahme beruht, das Polynom $Y$ lasse sich in lineare Faktoren zerlegen. Gerade hier, wo die allgemeine Möglichkeit dieser Zerlegung bewiesen werden soll, wäre diese Annahme nichts anderes als eine Vorwegnahme des zu Beweisenden. Dennoch haben sich nicht alle, die analytische Beweise des Hauptsatzes versucht haben, vor genau solchen Fehlschlüssen gehütet. Den Ursprung dieses täuschenden Anscheins haben wir schon in der Formulierung der Untersuchung selbst bemerkt: Alle haben nur nach der \textit{Form} der Wurzeln von Gleichungen gefragt, während sie deren voreilig angenommene \textit{Existenz} hätten beweisen müssen. Doch ist über dieses Vorgehen, das sich allzu weit von Strenge und Klarheit entfernt, schon in der oben angeführten Abhandlung genug gesagt worden. Wir werden daher nun die Sätze des vorigen Artikels,
\SourceLeaf{048}
von denen wir für unseren Zweck wenigstens einen nicht entbehren können, auf eine festere Grundlage stellen. Wir beginnen mit dem zweiten als dem leichteren.

\Needspace{18\baselineskip}\SecondArticle{8}
Wir bezeichnen mit $\rho$ die Funktion
\[
\begin{aligned}
&\frac{\pi(x-b)(x-c)(x-d)\, .\, .\, .\, .\, .}{(a-b)^2(a-c)^2(a-d)^2\, .\, .\, .\, .\, .}\\
&+\frac{\pi(x-a)(x-c)(x-d)\, .\, .\, .\, .\, .}{(b-a)^2(b-c)^2(b-d)^2\, .\, .\, .\, .\, .}\\
&+\frac{\pi(x-a)(x-b)(x-d)\, .\, .\, .\, .\, .}{(c-a)^2(c-b)^2(c-d)^2\, .\, .\, .\, .\, .}\\
&+\frac{\pi(x-a)(x-b)(x-c)\, .\, .\, .\, .\, .}{(d-a)^2(d-b)^2(d-c)^2\, .\, .\, .\, .\, .}\\
&+\text{usw.}
\end{aligned}
\]
die, weil $\pi$ durch jeden einzelnen Nenner teilbar ist, ein Polynom in den Unbestimmten $x,a,b,c$ usw. ist. Ferner setzen wir $\frac{\mathrm d\upsilon}{\mathrm dx}=\upsilon'$, sodass gilt
\[
\begin{aligned}
\upsilon'={}&(x-b)(x-c)(x-d)\, .\, .\, .\, .\, .\\
&+(x-a)(x-c)(x-d)\, .\, .\, .\, .\, .\\
&+(x-a)(x-b)(x-d)\, .\, .\, .\, .\, .\\
&+(x-a)(x-b)(x-c)\, .\, .\, .\, .\, .\\
&+\text{usw.}
\end{aligned}
\]
Für $x=a$ gilt offenbar $\rho\upsilon'=\pi$. Daraus schließen wir, dass das Polynom $\pi-\rho\upsilon'$ identisch durch $x-a$ und ebenso durch $x-b,x-c$ usw. sowie durch ihr Produkt $\upsilon$ teilbar ist. Setzen wir also
\[
\frac{\pi-\rho\upsilon'}{\upsilon}=\sigma
\]
so ist $\sigma$ ein Polynom in den Unbestimmten $x,a,b,c$ usw., das wie $\rho$ hinsichtlich der Unbestimmten $a,b,c$ usw. symmetrisch ist. Es lassen sich daher zwei Polynome $r,s$ in den Unbestimmten $x,l',l'',l'''$ usw. finden, die durch die Substitutionen $l'=\lambda',l''=\lambda'',l'''=\lambda'''$ usw. jeweils in $\rho,\sigma$ übergehen. Bezeichnen wir nun entsprechend das Polynom
\[
\begin{gathered}mx^{m-1}-(m-1)l'x^{m-2}+(m-2)l''x^{m-3}\\-(m-3)l'''x^{m-4}+\text{usw.}\end{gathered}
\]
also die Ableitung $\frac{\mathrm dy}{\mathrm dx}$, mit $y'$, sodass $y'$ durch dieselben
\SourceLeaf{049}
Substitutionen in $\upsilon'$ übergeht, so geht offenbar $p-sy-ry'$ durch dieselben Substitutionen in $\pi-\sigma\upsilon-\rho\upsilon'$, also in 0 über und muss daher schon für sich identisch verschwinden (Artikel 5). Wir haben somit die identische Gleichung
\[
p=sy+ry'
\]
Nehmen wir nun an, durch die Substitution $l'=L',l''=L'',l'''=L'''$ usw. ergebe sich $r=R,s=S$, so gilt ebenfalls identisch
\[
P=SY+RY'
\]
Da hier $S,R$ Polynome in $x$ sind, $P$ dagegen eine bestimmte Größe oder Zahl ist, folgt von selbst, dass $Y$ und $Y'$ keinen gemeinsamen Teiler haben können, wenn nicht $P=0$ ist. Dies ist genau der zweite Satz von Artikel 6.

